\documentclass[11pt,a4paper]{article}

\usepackage[a4paper,margin=25mm]{geometry}
\usepackage{fontspec}
\usepackage{polyglossia}
\setmainlanguage{english}
\setmainfont{DejaVu Serif}
\setsansfont{DejaVu Sans}
\setmonofont{DejaVu Sans Mono}

\usepackage{amsmath,amssymb,amsfonts,mathtools}
\usepackage{booktabs,longtable,array,tabularx}
\usepackage{enumitem}
\usepackage{hyperref}
\usepackage{xcolor}
\usepackage{listings}
\usepackage{fancyhdr}
\usepackage{titlesec}
\usepackage{parskip}

\hypersetup{
    colorlinks=true,
    linkcolor=blue!60!black,
    urlcolor=blue!60!black,
    citecolor=blue!60!black,
    pdftitle={Dnipro User's Guide by Examples},
    pdfauthor={dnipro project}
}

\setlength{\headheight}{14pt}
\pagestyle{fancy}
\fancyhf{}
\lhead{\texttt{dnipro} User's Guide}
\rhead{v0.1.0}
\cfoot{\thepage}

\lstdefinelanguage{spice}{
    morekeywords={DC,AC,SIN,PULSE,PWL,SQUARE,SAW,op,run,print,quit},
    morecomment=[l]{*},
    morecomment=[l]{.},
    morestring=[b]",
    sensitive=false
}

\lstdefinestyle{basecode}{
    basicstyle=\ttfamily\small,
    breaklines=true,
    frame=single,
    columns=fullflexible,
    keepspaces=true,
    showstringspaces=false,
    backgroundcolor=\color{gray!5},
    rulecolor=\color{gray!40}
}

\lstdefinestyle{python}{
    style=basecode,
    language=Python
}

\lstdefinestyle{spice}{
    style=basecode,
    language=spice
}

\lstdefinestyle{shell}{
    style=basecode,
    backgroundcolor=\color{black!4},
    rulecolor=\color{gray!60}
}

\lstdefinestyle{textout}{
    style=basecode,
    backgroundcolor=\color{blue!3},
    rulecolor=\color{blue!30}
}

\lstset{style=basecode}

\titleformat{\section}{\Large\bfseries}{\thesection}{1em}{}
\titleformat{\subsection}{\large\bfseries}{\thesubsection}{1em}{}
\titleformat{\subsubsection}{\normalsize\bfseries}{\thesubsubsection}{1em}{}

\newcommand{\project}{\texttt{dnipro}}
\newcommand{\code}[1]{\texttt{#1}}
\newcommand{\sympy}{\texttt{SymPy}}
\newcommand{\ngspice}{\texttt{ngspice}}
\newcommand{\mna}{Modified Nodal Analysis}
\newcommand{\examplesroot}{\path{../examples/user_guide}}

\begin{document}

\begin{titlepage}
    \centering
    \vspace*{20mm}
    {\Huge\bfseries Dnipro User's Guide\par}
    \vspace{6mm}
    {\Large\bfseries By Examples\par}
    \vspace{12mm}
    {\large Symbolic equation-based MNA for SPICE-like circuits\par}
    \vspace{18mm}
    {\large Version 0.1.0\par}
    \vfill
    \begin{tabular}{ll}
        Package name: & \project \\
        Implementation language: & Python 3.11+ \\
        Symbolic engine: & \sympy \\
        Optional reference simulator: & \ngspice \\
        Example root: & \examplesroot \\
    \end{tabular}
    \vfill
    {\large \today\par}
\end{titlepage}

\tableofcontents
\newpage

\section{Purpose and Scope}

This guide demonstrates \project{} through runnable examples. Each example is
backed by files in \path{examples/user_guide}: a \project{} netlist, a runnable
\ngspice{} template, and a Python script that exercises the current public API.

The guide is intentionally example-driven. It is not a replacement for the
reference manual. The reference manual describes the implemented API surface;
this document shows how that API behaves on concrete circuits and how the
resulting symbolic systems map to standard \mna{} notation.

\subsection{What the examples demonstrate}

\begin{itemize}[noitemsep]
    \item Residual equation construction \(F(x,\dot{x},t,p)=0\).
    \item Classical linear extraction \(A x = b\) where the residuals are linear
          in the unknown vector.
    \item Time-domain descriptor extraction \(Gx + C\dot{x} = b\).
    \item Laplace-domain linear systems \(A(s)x=b\).
    \item Numeric and frequency substitution using exact \sympy{} numbers.
    \item Nonlinear residual construction for the diode model.
    \item Expected API errors for unsupported transformations.
\end{itemize}

\subsection{Repository layout}

Every example directory uses the same structure:

\begin{lstlisting}[style=textout]
examples/user_guide/NN_example_name/
  dnipro.cir    # dnipro-compatible symbolic netlist
  ngspice.cir   # runnable numeric ngspice template
  run.py        # Python script using the dnipro API
\end{lstlisting}

Run a Python example from the repository root:

\begin{lstlisting}[style=shell]
/home/konst/Documents/Projects/venv/bin/python examples/user_guide/01_voltage_divider/run.py
\end{lstlisting}

Run the matching \ngspice{} template:

\begin{lstlisting}[style=shell]
ngspice -b examples/user_guide/01_voltage_divider/ngspice.cir
\end{lstlisting}

\section{Conventions}

\subsection{SPICE-compatible directions}

\project{} follows the SPICE convention for voltage and current directions.
For a voltage source

\begin{lstlisting}[style=spice]
V1 nplus nminus value
\end{lstlisting}

the source voltage is \(v(nplus)-v(nminus)\). The branch-current unknown
\(i_{V1}\) follows the current direction through the source from the first node
to the second node.

For a current source

\begin{lstlisting}[style=spice]
I1 nplus nminus value
\end{lstlisting}

the source current is defined from the first node to the second node. In KCL
residuals, \project{} writes currents leaving each node as positive. Therefore a
source \code{I1 out 0 Iin} contributes \(+I_{in}\) to the KCL residual at
\code{out}.

\subsection{Unknown ordering}

The unknown vector always lists node voltages first, in parser insertion order,
followed by branch currents in first-appearance order. In time domain, unknowns
are \sympy{} functions such as \(v_{out}(t)\). In Laplace domain, unknowns are
plain symbols such as \(v_{out}\).

\subsection{Residual, linear, and descriptor forms}

The primary representation is the residual system

\begin{equation}
    F(x,\dot{x},t,p)=0.
\end{equation}

If the residuals are algebraically linear in \(x\), \project{} can derive

\begin{equation}
    A x = b.
\end{equation}

For linear time-domain systems with derivatives, \project{} can derive the
descriptor form

\begin{equation}
    Gx + C\dot{x} = b.
\end{equation}

Classical \(A x=b\) extraction is not the right time-domain view when residuals
contain derivatives such as \(\frac{d}{dt}v_{out}(t)\). Use descriptor form in
time domain, or use Laplace domain to obtain \(A(s)x=b\).

\section{Shared Python Helpers}

All example scripts import the shared helper module:

\lstinputlisting[style=python]{../examples/user_guide/common.py}

The helpers are deliberately thin wrappers around the public \project{} API:

\begin{itemize}[noitemsep]
    \item \code{parse\_netlist\_file};
    \item \code{build\_equations};
    \item \code{extract\_linear\_system};
    \item \code{extract\_descriptor\_system};
    \item the text exporters.
\end{itemize}

\section{Example 1: Voltage Divider}
\label{sec:ex-voltage-divider}

This example is the smallest useful linear circuit. It demonstrates resistor
KCL, voltage-source branch current allocation, residual reporting, and
classical \(A x=b\) extraction.

\subsection{dnipro netlist}

\lstinputlisting[style=spice]{../examples/user_guide/01_voltage_divider/dnipro.cir}

\subsection{ngspice template}

\lstinputlisting[style=spice]{../examples/user_guide/01_voltage_divider/ngspice.cir}

\subsection{Python script}

\lstinputlisting[style=python]{../examples/user_guide/01_voltage_divider/run.py}

\subsection{Symbolic system}

The unknown vector is

\begin{equation}
    x =
    \begin{bmatrix}
        v_{in}(t) \\
        v_{out}(t) \\
        i_{V1}(t)
    \end{bmatrix}.
\end{equation}

The residual equations are

\begin{align}
    \frac{v_{in}-v_{out}}{R_1} + i_{V1} &= 0, \\
    \frac{v_{out}-v_{in}}{R_1} + \frac{v_{out}}{R_2} &= 0, \\
    v_{in} - V_{in} &= 0.
\end{align}

Since the system is linear in \(x\), \project{} extracts

\begin{equation}
    A =
    \begin{bmatrix}
        \frac{1}{R_1} & -\frac{1}{R_1} & 1 \\
        -\frac{1}{R_1} & \frac{1}{R_1}+\frac{1}{R_2} & 0 \\
        1 & 0 & 0
    \end{bmatrix},
    \quad
    b =
    \begin{bmatrix}
        0 \\
        0 \\
        V_{in}
    \end{bmatrix}.
\end{equation}

With \(V_{in}=10\), \(R_1=R_2=1\,\mathrm{k}\Omega\), the \ngspice{} template
reports \(v(out)=5\,\mathrm{V}\). The sign of \(i(v1)\) follows \ngspice{}'s
source-current convention.

\section{Example 2: Current Source Direction}
\label{sec:ex-current-source-direction}

This example isolates the SPICE current-source direction convention.

\subsection{dnipro netlist}

\lstinputlisting[style=spice]{../examples/user_guide/02_current_source_direction/dnipro.cir}

\subsection{ngspice template}

\lstinputlisting[style=spice]{../examples/user_guide/02_current_source_direction/ngspice.cir}

\subsection{Python script}

\lstinputlisting[style=python]{../examples/user_guide/02_current_source_direction/run.py}

\subsection{Symbolic system}

The source \code{I1 out 0 Iin} defines current from \code{out} to ground.
Because \project{} writes KCL as a sum of currents leaving the node, the residual
is

\begin{equation}
    I_{in} + \frac{v_{out}}{R} = 0.
\end{equation}

The extracted scalar system is

\begin{equation}
    A = \begin{bmatrix}\frac{1}{R}\end{bmatrix},
    \quad
    x = \begin{bmatrix}v_{out}(t)\end{bmatrix},
    \quad
    b = \begin{bmatrix}-I_{in}\end{bmatrix}.
\end{equation}

For \(I_{in}=1\,\mathrm{mA}\) and \(R=1\,\mathrm{k}\Omega\), the operating point
is \(v(out)=-1\,\mathrm{V}\). This negative voltage is expected: a positive
current source from \code{out} to ground removes current from the node.

\section{Example 3: RC Time-Domain Descriptor System}
\label{sec:ex-rc-time-descriptor}

This example introduces a capacitor and therefore a time derivative. The primary
time-domain report is the residual system, and the matrix view is the descriptor
form \(Gx + C\dot{x}=b\).

\subsection{dnipro netlist}

\lstinputlisting[style=spice]{../examples/user_guide/03_rc_time_descriptor/dnipro.cir}

\subsection{ngspice template}

\lstinputlisting[style=spice]{../examples/user_guide/03_rc_time_descriptor/ngspice.cir}

\subsection{Python script}

\lstinputlisting[style=python]{../examples/user_guide/03_rc_time_descriptor/run.py}

\subsection{Residual equations}

The unknown vector is

\begin{equation}
    x =
    \begin{bmatrix}
        v_{in}(t) \\
        v_{out}(t) \\
        i_{V1}(t)
    \end{bmatrix},
    \quad
    \dot{x} =
    \begin{bmatrix}
        \dot{v}_{in}(t) \\
        \dot{v}_{out}(t) \\
        \dot{i}_{V1}(t)
    \end{bmatrix}.
\end{equation}

The residual equations are

\begin{align}
    \frac{v_{in}-v_{out}}{R} + i_{V1} &= 0, \\
    \frac{v_{out}-v_{in}}{R} + C\frac{d v_{out}}{dt} &= 0, \\
    v_{in} - V_{in} &= 0.
\end{align}

\subsection{Descriptor matrix}

\project{} extracts

\begin{equation}
    G =
    \begin{bmatrix}
        \frac{1}{R} & -\frac{1}{R} & 1 \\
        -\frac{1}{R} & \frac{1}{R} & 0 \\
        1 & 0 & 0
    \end{bmatrix},
    \quad
    C =
    \begin{bmatrix}
        0 & 0 & 0 \\
        0 & C & 0 \\
        0 & 0 & 0
    \end{bmatrix},
    \quad
    b =
    \begin{bmatrix}
        0 \\
        0 \\
        V_{in}
    \end{bmatrix}.
\end{equation}

The nonzero \(C_{2,2}\) entry comes from the capacitor current
\(C\,d v_{out}/dt\). Classical \(A x=b\) extraction is not used here because the
time-domain residuals contain derivatives.

\section{Example 4: RC Low-Pass in Laplace Domain}
\label{sec:ex-rc-laplace}

The same RC circuit can be built in Laplace domain. In that mode the capacitor
contribution becomes the algebraic admittance \(sC\), and the circuit has a
classical linear system \(A(s)x=b\).

\subsection{dnipro netlist}

\lstinputlisting[style=spice]{../examples/user_guide/04_rc_laplace/dnipro.cir}

\subsection{ngspice template}

\lstinputlisting[style=spice]{../examples/user_guide/04_rc_laplace/ngspice.cir}

\subsection{Python script}

\lstinputlisting[style=python]{../examples/user_guide/04_rc_laplace/run.py}

\subsection{Laplace-domain matrix}

The unknown vector is

\begin{equation}
    x =
    \begin{bmatrix}
        v_{in} \\
        v_{out} \\
        i_{V1}
    \end{bmatrix}.
\end{equation}

The extracted matrix is

\begin{equation}
    A(s) =
    \begin{bmatrix}
        \frac{1}{R} & -\frac{1}{R} & 1 \\
        -\frac{1}{R} & \frac{1}{R}+sC & 0 \\
        1 & 0 & 0
    \end{bmatrix},
    \quad
    b =
    \begin{bmatrix}
        0 \\
        0 \\
        V_{in}
    \end{bmatrix}.
\end{equation}

The example script also demonstrates frequency substitution:

\begin{align}
    s &= j\omega, \\
    s &= j2\pi f.
\end{align}

For \code{--freq 50}, \project{} records \(s = 100j\pi\) in the residual system
metadata.

\section{Example 5: Nonlinear Diode Residual}
\label{sec:ex-diode-nonlinear}

This example demonstrates the nonlinear diode component. The circuit is
symbolically supported as a residual system, but it is intentionally rejected by
classical linear extraction.

\subsection{dnipro netlist}

\lstinputlisting[style=spice]{../examples/user_guide/08_diode_nonlinear/dnipro.cir}

\subsection{ngspice template}

\lstinputlisting[style=spice]{../examples/user_guide/08_diode_nonlinear/ngspice.cir}

\subsection{Python script}

\lstinputlisting[style=python]{../examples/user_guide/08_diode_nonlinear/run.py}

\subsection{Nonlinear residual}

The unknown vector is

\begin{equation}
    x =
    \begin{bmatrix}
        v_{in}(t) \\
        v_{out}(t) \\
        i_{V1}(t)
    \end{bmatrix}.
\end{equation}

The diode current is

\begin{equation}
    i_D =
    I_s\left(\exp\!\left(\frac{v_{out}}{nV_t}\right)-1\right).
\end{equation}

The residual equations are

\begin{align}
    \frac{v_{in}-v_{out}}{R} + i_{V1} &= 0, \\
    \frac{v_{out}-v_{in}}{R}
        + I_s\left(\exp\!\left(\frac{v_{out}}{nV_t}\right)-1\right) &= 0, \\
    v_{in} - V_{in} &= 0.
\end{align}

\subsection{Expected linear extraction error}

This residual system is nonlinear in \(v_{out}\). Calling
\code{extract\_linear\_system} therefore raises \code{NonlinearSystemError}.
This is expected behavior: \project{} has built the correct symbolic nonlinear
system, but it does not pretend that the system can be represented as a
classical linear \(A x=b\) matrix.

The example script also applies exact numeric substitutions:

\begin{lstlisting}[style=textout]
Vin=5, R=1k, Is=1e-14, n=1, Vt=25.85m
\end{lstlisting}

These values are parsed into exact \sympy{} numbers before substitution.

\section{Example 6: Voltage Source Direction}
\label{sec:ex-voltage-source-direction}

The previous examples used a voltage source from a named node to ground. This
example reverses the source terminals and demonstrates the resulting sign in
the voltage constraint.

\subsection{dnipro netlist}

\lstinputlisting[style=spice]{../examples/user_guide/09_voltage_source_direction/dnipro.cir}

\subsection{ngspice template}

\lstinputlisting[style=spice]{../examples/user_guide/09_voltage_source_direction/ngspice.cir}

\subsection{Python script}

\lstinputlisting[style=python]{../examples/user_guide/09_voltage_source_direction/run.py}

\subsection{Sign convention}

For \code{V1 0 out Vin}, the source voltage is \(v(0)-v(out)\). Since ground is
zero, the voltage constraint is

\begin{equation}
    -v_{out} - V_{in} = 0.
\end{equation}

The KCL equation at \code{out} is

\begin{equation}
    \frac{v_{out}}{R} - i_{V1} = 0,
\end{equation}

because the branch current \(i_{V1}\) is defined from ground to \code{out}. With
\(V_{in}=5\), \ngspice{} reports \(v(out)=-5\,\mathrm{V}\), which is exactly the
source orientation encoded by the netlist.

\section{Example 7: Floating Capacitor}
\label{sec:ex-floating-capacitor}

This example places the capacitor between two non-ground nodes. It demonstrates
that the capacitor derivative contributes to both KCL equations with opposite
signs.

\subsection{dnipro netlist}

\lstinputlisting[style=spice]{../examples/user_guide/10_floating_capacitor/dnipro.cir}

\subsection{ngspice template}

\lstinputlisting[style=spice]{../examples/user_guide/10_floating_capacitor/ngspice.cir}

\subsection{Python script}

\lstinputlisting[style=python]{../examples/user_guide/10_floating_capacitor/run.py}

\subsection{Residual structure}

The key capacitor current is

\begin{equation}
    i_C = C\frac{d}{dt}\left(v_{in}-v_{out}\right).
\end{equation}

Therefore the two KCL residuals contain \(+i_C\) at \code{in} and \(-i_C\) at
\code{out}. This is a useful conservation check for dynamic two-terminal
elements.

\section{Example 8: RL Branch-Current Formulation}
\label{sec:ex-rl-branch-current}

Inductors are represented with branch-current unknowns in both time and Laplace
domains. This keeps the formulation aligned with SPICE-style MNA.

\subsection{dnipro netlist}

\lstinputlisting[style=spice]{../examples/user_guide/05_rl_branch_current/dnipro.cir}

\subsection{ngspice template}

\lstinputlisting[style=spice]{../examples/user_guide/05_rl_branch_current/ngspice.cir}

\subsection{Python script}

\lstinputlisting[style=python]{../examples/user_guide/05_rl_branch_current/run.py}

\subsection{Inductor equation}

The time-domain branch equation is

\begin{equation}
    v_{in} - v_{out} - L\frac{d i_{L1}}{dt} = 0.
\end{equation}

In Laplace domain the same relation becomes

\begin{equation}
    v_{in} - v_{out} - sL\,i_{L1} = 0.
\end{equation}

The branch current \(i_{L1}\) is part of the unknown vector. This is the same
principle used for voltage sources and controlled voltage sources.

\section{Example 9: RL Descriptor Form}
\label{sec:ex-rl-descriptor}

This variant connects the inductor from \code{out} to ground and shows how the
inductor derivative appears in the descriptor matrix.

\subsection{dnipro netlist}

\lstinputlisting[style=spice]{../examples/user_guide/11_rl_descriptor/dnipro.cir}

\subsection{ngspice template}

\lstinputlisting[style=spice]{../examples/user_guide/11_rl_descriptor/ngspice.cir}

\subsection{Python script}

\lstinputlisting[style=python]{../examples/user_guide/11_rl_descriptor/run.py}

\subsection{Descriptor contribution}

For the unknown vector

\begin{equation}
    x =
    \begin{bmatrix}
        v_{in}(t) \\
        v_{out}(t) \\
        i_{V1}(t) \\
        i_{L1}(t)
    \end{bmatrix},
\end{equation}

the inductor equation contributes a dynamic coefficient \(-L\) in the row for
\code{L1} and the column for \(\dot{i}_{L1}\). The descriptor form is therefore
the natural time-domain matrix representation for this circuit.

\section{Example 10: RLC Low-Pass}
\label{sec:ex-rlc-lowpass}

This example combines a resistor, an inductor branch current, and a capacitor
voltage derivative in one circuit.

\subsection{dnipro netlist}

\lstinputlisting[style=spice]{../examples/user_guide/06_rlc_lowpass/dnipro.cir}

\subsection{ngspice template}

\lstinputlisting[style=spice]{../examples/user_guide/06_rlc_lowpass/ngspice.cir}

\subsection{Python script}

\lstinputlisting[style=python]{../examples/user_guide/06_rlc_lowpass/run.py}

\subsection{Mixed dynamic structure}

The residual system contains both dynamic patterns:

\begin{itemize}[noitemsep]
    \item the capacitor contributes \(C\,d v_{out}/dt\) to the KCL equation at
          \code{out};
    \item the inductor contributes the branch equation
          \(v_{n1}-v_{out}-L\,d i_{L1}/dt=0\).
\end{itemize}

In time domain, use descriptor extraction. In Laplace domain, the same example
has a classical \(A(s)x=b\) matrix.

\section{Example 11: Sinusoidal Voltage Source}
\label{sec:ex-sine-voltage}

This example shows a symbolic sinusoidal voltage source in time domain.

\subsection{dnipro netlist}

\lstinputlisting[style=spice]{../examples/user_guide/12_sine_voltage_source/dnipro.cir}

\subsection{ngspice template}

\lstinputlisting[style=spice]{../examples/user_guide/12_sine_voltage_source/ngspice.cir}

\subsection{Python script}

\lstinputlisting[style=python]{../examples/user_guide/12_sine_voltage_source/run.py}

\subsection{Source expression}

The source

\begin{lstlisting}[style=spice]
SIN(0 Vm omega)
\end{lstlisting}

is represented as

\begin{equation}
    V_m\sin(\omega t).
\end{equation}

The voltage-source residual is therefore

\begin{equation}
    v_{in}(t) - V_m\sin(\omega t) = 0.
\end{equation}

\section{Example 12: Sinusoidal Current Source}
\label{sec:ex-sine-current}

The current-source version demonstrates that the same source expression is
inserted into KCL using SPICE current direction.

\subsection{dnipro netlist}

\lstinputlisting[style=spice]{../examples/user_guide/13_sine_current_source/dnipro.cir}

\subsection{ngspice template}

\lstinputlisting[style=spice]{../examples/user_guide/13_sine_current_source/ngspice.cir}

\subsection{Python script}

\lstinputlisting[style=python]{../examples/user_guide/13_sine_current_source/run.py}

\subsection{KCL contribution}

For \code{I1 out 0 SIN(0 Im omega)}, the KCL residual at \code{out} includes

\begin{equation}
    I_m\sin(\omega t) + \frac{v_{out}}{R}
        + C\frac{d v_{out}}{dt} = 0.
\end{equation}

The positive sign on the source term follows directly from the current direction
from \code{out} to ground.

\section{Example 13: PULSE Source}
\label{sec:ex-pulse}

The PULSE source is expanded into a symbolic \sympy{} \code{Piecewise}
expression.

\subsection{dnipro netlist}

\lstinputlisting[style=spice]{../examples/user_guide/14_pulse_source/dnipro.cir}

\subsection{ngspice template}

\lstinputlisting[style=spice]{../examples/user_guide/14_pulse_source/ngspice.cir}

\subsection{Python script}

\lstinputlisting[style=python]{../examples/user_guide/14_pulse_source/run.py}

\subsection{Timing parameters}

The example uses the supported seven-argument source form:

\begin{lstlisting}[style=spice]
PULSE(0 Vh 0 tr tf ton T)
\end{lstlisting}

The duty-related timing is explicit: rise time, fall time, on time, and period
are all supplied. \project{} keeps those symbols exact and builds the resulting
piecewise source expression.

\section{Example 14: PWL Source}
\label{sec:ex-pwl}

Piecewise-linear sources are also converted to \sympy{} \code{Piecewise}
expressions.

\subsection{dnipro netlist}

\lstinputlisting[style=spice]{../examples/user_guide/15_pwl_source/dnipro.cir}

\subsection{ngspice template}

\lstinputlisting[style=spice]{../examples/user_guide/15_pwl_source/ngspice.cir}

\subsection{Python script}

\lstinputlisting[style=python]{../examples/user_guide/15_pwl_source/run.py}

\subsection{Exact time points}

The source

\begin{lstlisting}[style=spice]
PWL(0 0 1m 5 2m 0)
\end{lstlisting}

uses exact rational time values internally: \(1\,\mathrm{ms}=1/1000\) and
\(2\,\mathrm{ms}=1/500\). This is important for stable symbolic output and JSON
serialization.

\section{Example 15: Dependent Sources}
\label{sec:ex-dependent-sources}

This example combines all four dependent source families in one circuit:
VCVS, VCCS, CCCS, and CCVS.

\subsection{dnipro netlist}

\lstinputlisting[style=spice]{../examples/user_guide/07_dependent_sources/dnipro.cir}

\subsection{ngspice template}

\lstinputlisting[style=spice]{../examples/user_guide/07_dependent_sources/ngspice.cir}

\subsection{Python script}

\lstinputlisting[style=python]{../examples/user_guide/07_dependent_sources/run.py}

\subsection{Device roles}

\begin{longtable}{p{0.16\textwidth}p{0.24\textwidth}p{0.50\textwidth}}
\toprule
Type & Example & Role \\
\midrule
VCVS & \code{E1} & Adds its own branch current and a controlled voltage
constraint. \\
VCCS & \code{G1} & Adds a controlled current contribution to KCL without a new
branch current. \\
CCCS & \code{F1} & Uses the existing current through \code{Vsense} as the control
current. \\
CCVS & \code{H1} & Adds its own branch current and uses \code{Vsense} as the
control current. \\
\bottomrule
\end{longtable}

Current-controlled sources require a control branch. The examples use a
zero-volt voltage source \code{Vsense}, which is the standard SPICE technique
for measuring a branch current.

\section{Example 16: Parser RVI Details}
\label{sec:ex-parser-rvi}

This example demonstrates parser behavior for full-line comments, ground aliases,
numeric values, and \code{DC(...)} source syntax.

\subsection{dnipro netlist}

\lstinputlisting[style=spice]{../examples/user_guide/16_parser_rvi_comments_ground_alias/dnipro.cir}

\subsection{ngspice template}

\lstinputlisting[style=spice]{../examples/user_guide/16_parser_rvi_comments_ground_alias/ngspice.cir}

\subsection{Python script}

\lstinputlisting[style=python]{../examples/user_guide/16_parser_rvi_comments_ground_alias/run.py}

\subsection{Notes}

The parser ignores full-line comments beginning with \code{*}. Ground aliases
such as \code{gnd} normalize to the internal ground node. The resistor value
\code{1k} is parsed as the exact integer \(1000\), while \code{DC(Iload)} is
reduced to the symbolic source value \(Iload\).

\section{Example 17: Parser C/L Values}
\label{sec:ex-parser-cl}

This parser-oriented example checks capacitor and inductor lines.

\subsection{dnipro netlist}

\lstinputlisting[style=spice]{../examples/user_guide/17_parser_cl_values/dnipro.cir}

\subsection{ngspice template}

\lstinputlisting[style=spice]{../examples/user_guide/17_parser_cl_values/ngspice.cir}

\subsection{Python script}

\lstinputlisting[style=python]{../examples/user_guide/17_parser_cl_values/run.py}

\subsection{Notes}

The \project{} netlist is intentionally minimal because it mirrors the parser
unit test. The \ngspice{} wrapper adds an excitation source and a large leak
resistor so that the external simulator has a nonsingular operating point before
AC analysis. The symbolic \project{} example does not need that wrapper element.

\section{Example 18: Numeric Diode Model Values}
\label{sec:ex-diode-model-values}

This example focuses on diode model parsing rather than circuit solving.

\subsection{dnipro netlist}

\lstinputlisting[style=spice]{../examples/user_guide/18_diode_model_values/dnipro.cir}

\subsection{ngspice template}

\lstinputlisting[style=spice]{../examples/user_guide/18_diode_model_values/ngspice.cir}

\subsection{Python script}

\lstinputlisting[style=python]{../examples/user_guide/18_diode_model_values/run.py}

\subsection{Exact values}

The parser converts the model

\begin{lstlisting}[style=spice]
.model Ddefault D(Is=1e-14 n=1 Vt=25.85m)
\end{lstlisting}

into exact \sympy{} numbers:

\begin{equation}
    I_s = \frac{1}{10^{14}}, \quad n = 1, \quad V_t = \frac{517}{20000}.
\end{equation}

\section{Example 19: Dependent Source Parser Bundle}
\label{sec:ex-dependent-parser-bundle}

This parser bundle contains all four dependent source types without the
additional elements needed for a complete circuit solve.

\subsection{dnipro netlist}

\lstinputlisting[style=spice]{../examples/user_guide/19_dependent_source_parser_bundle/dnipro.cir}

\subsection{ngspice template}

\lstinputlisting[style=spice]{../examples/user_guide/19_dependent_source_parser_bundle/ngspice.cir}

\subsection{Python script}

\lstinputlisting[style=python]{../examples/user_guide/19_dependent_source_parser_bundle/run.py}

\subsection{Notes}

The \project{} netlist demonstrates typed parsing only. The \ngspice{} wrapper
adds independent sources and loads so the dependent devices can be evaluated in a
runnable operating-point analysis.

\section{Example 20: Symbolic Diode Model Parser}
\label{sec:ex-diode-parser-symbolic}

This example contains a diode and symbolic diode model without the surrounding
source and resistor used in the full nonlinear example.

\subsection{dnipro netlist}

\lstinputlisting[style=spice]{../examples/user_guide/20_diode_parser_symbolic_model/dnipro.cir}

\subsection{ngspice template}

\lstinputlisting[style=spice]{../examples/user_guide/20_diode_parser_symbolic_model/ngspice.cir}

\subsection{Python script}

\lstinputlisting[style=python]{../examples/user_guide/20_diode_parser_symbolic_model/run.py}

\subsection{Residual}

With only the diode connected from \code{out} to ground, the KCL residual is the
diode current itself:

\begin{equation}
    I_s\left(\exp\!\left(\frac{v_{out}}{nV_t}\right)-1\right)=0.
\end{equation}

\section{Example 21: Single-Capacitor Frequency Substitution}
\label{sec:ex-single-cap-frequency}

This example isolates the Laplace-domain capacitor term and frequency
substitution machinery.

\subsection{dnipro netlist}

\lstinputlisting[style=spice]{../examples/user_guide/21_single_capacitor_frequency/dnipro.cir}

\subsection{ngspice template}

\lstinputlisting[style=spice]{../examples/user_guide/21_single_capacitor_frequency/ngspice.cir}

\subsection{Python script}

\lstinputlisting[style=python]{../examples/user_guide/21_single_capacitor_frequency/run.py}

\subsection{Frequency substitution}

In Laplace domain the residual is

\begin{equation}
    sC\,v_{out} = 0.
\end{equation}

For \(f=50\,\mathrm{Hz}\), \project{} substitutes

\begin{equation}
    s = j2\pi f = 100j\pi.
\end{equation}

The \ngspice{} wrapper adds a current excitation and a large leak resistor only
to make the AC analysis runnable.

\section{Example 22: Missing Numeric Parameters}
\label{sec:ex-missing-params}

Numeric substitution is strict: every symbolic parameter present in the residual
system must be provided.

\subsection{dnipro netlist}

\lstinputlisting[style=spice]{../examples/user_guide/22_numeric_substitution_missing_params/dnipro.cir}

\subsection{ngspice template}

\lstinputlisting[style=spice]{../examples/user_guide/22_numeric_substitution_missing_params/ngspice.cir}

\subsection{Python script}

\lstinputlisting[style=python]{../examples/user_guide/22_numeric_substitution_missing_params/run.py}

\subsection{Expected error}

The script intentionally calls \code{apply\_numeric\_substitutions} with no
parameters. The expected result is \code{MissingParameterError}, reporting
\code{C}, \code{R}, and \code{Vin}.

\section{Example 23: Linear Extraction Rejects Diode}
\label{sec:ex-linear-rejects-diode}

This example repeats the diode circuit but focuses only on the expected failure
mode of classical linear extraction.

\subsection{dnipro netlist}

\lstinputlisting[style=spice]{../examples/user_guide/23_linear_extraction_rejects_diode/dnipro.cir}

\subsection{ngspice template}

\lstinputlisting[style=spice]{../examples/user_guide/23_linear_extraction_rejects_diode/ngspice.cir}

\subsection{Python script}

\lstinputlisting[style=python]{../examples/user_guide/23_linear_extraction_rejects_diode/run.py}

\subsection{Expected error}

The residual system is valid but nonlinear. Calling
\code{extract\_linear\_system} raises \code{NonlinearSystemError}. This explicit
failure prevents accidental use of a linear matrix representation for a
nonlinear circuit.

\section{Example 24: Descriptor Rejects Laplace Domain}
\label{sec:ex-descriptor-rejects-laplace}

Descriptor extraction is defined for time-domain systems only.

\subsection{dnipro netlist}

\lstinputlisting[style=spice]{../examples/user_guide/24_descriptor_rejects_laplace/dnipro.cir}

\subsection{ngspice template}

\lstinputlisting[style=spice]{../examples/user_guide/24_descriptor_rejects_laplace/ngspice.cir}

\subsection{Python script}

\lstinputlisting[style=python]{../examples/user_guide/24_descriptor_rejects_laplace/run.py}

\subsection{Expected error}

The Laplace-domain residual \(sC\,v_{out}=0\) is algebraic. Requesting
descriptor extraction from that system raises \code{NonlinearSystemError} with a
message explaining that descriptor form is available only for time-domain
systems.

\section{Example 25: Case-Preserved Component Name}
\label{sec:ex-case-preserved}

Component type dispatch is case-insensitive, but the original component name is
preserved.

\subsection{dnipro netlist}

\lstinputlisting[style=spice]{../examples/user_guide/25_case_preserved_component_name/dnipro.cir}

\subsection{ngspice template}

\lstinputlisting[style=spice]{../examples/user_guide/25_case_preserved_component_name/ngspice.cir}

\subsection{Python script}

\lstinputlisting[style=python]{../examples/user_guide/25_case_preserved_component_name/run.py}

\subsection{Notes}

The line \code{r\_load in out Rload} is parsed as a resistor, while the component
name remains \code{r\_load}. This matters for stable reporting and for preserving
the user's naming choices.

\section{Example 26: Expected Parser Errors}
\label{sec:ex-invalid-examples}

The final runnable example demonstrates unsupported or malformed syntax. The
\code{dnipro.cir} file is intentionally valid; the invalid cases are held inside
the Python script so that the example itself can run successfully.

\subsection{dnipro placeholder netlist}

\lstinputlisting[style=spice]{../examples/user_guide/26_invalid_examples_expected_errors/dnipro.cir}

\subsection{ngspice placeholder template}

\lstinputlisting[style=spice]{../examples/user_guide/26_invalid_examples_expected_errors/ngspice.cir}

\subsection{Python script}

\lstinputlisting[style=python]{../examples/user_guide/26_invalid_examples_expected_errors/run.py}

\subsection{Rejected syntax}

The rejected cases include invalid component names, unsupported arithmetic in
values, comma-separated source syntax, unsupported directives, malformed diode
lines, and incomplete diode model definitions. These are parser/API boundaries,
not simulator failures.

\appendix

\section{Complete Runnable Example Catalog}
\label{app:catalog}

All examples below are present in \code{examples/user\_guide}. Each directory
has the same three-file structure:

\begin{itemize}[noitemsep]
    \item \code{dnipro.cir};
    \item \code{ngspice.cir};
    \item \code{run.py}.
\end{itemize}

\begingroup
\scriptsize
\begin{longtable}{p{0.08\textwidth}p{0.38\textwidth}p{0.44\textwidth}}
\toprule
ID & Directory & Demonstrated feature \\
\midrule
\endfirsthead
\toprule
ID & Directory & Demonstrated feature \\
\midrule
\endhead
01 & \path{01_voltage_divider} & Basic resistor divider, residuals, and
classical \(A x=b\). \\
02 & \path{02_current_source_direction} & SPICE current-source direction and
sign of node voltage. \\
03 & \path{03_rc_time_descriptor} & RC time-domain residuals and
\(Gx+C\dot{x}=b\). \\
04 & \path{04_rc_laplace} & RC Laplace-domain \(A(s)x=b\) and frequency
substitution. \\
05 & \path{05_rl_branch_current} & Inductor branch-current formulation in
time and Laplace domains. \\
06 & \path{06_rlc_lowpass} & Mixed RLC dynamics with capacitor derivative and
inductor branch current. \\
07 & \path{07_dependent_sources} & VCVS, VCCS, CCCS, and CCVS in a complete
circuit. \\
08 & \path{08_diode_nonlinear} & Nonlinear diode residual and expected
linear-extraction rejection. \\
09 & \path{09_voltage_source_direction} & Reversed voltage-source terminals
and sign convention. \\
10 & \path{10_floating_capacitor} & Capacitor derivative contribution to two
non-ground nodes. \\
11 & \path{11_rl_descriptor} & Descriptor matrix for an inductor to ground. \\
12 & \path{12_sine_voltage_source} & Symbolic sinusoidal voltage source. \\
13 & \path{13_sine_current_source} & Symbolic sinusoidal current source and
KCL sign. \\
14 & \path{14_pulse_source} & PULSE source expansion to symbolic
\code{Piecewise}. \\
15 & \path{15_pwl_source} & PWL source expansion with exact rational times. \\
16 & \path{16_parser_rvi_comments_ground_alias} & Comments, ground aliases,
numeric values, and \code{DC(...)} syntax. \\
17 & \path{17_parser_cl_values} & Capacitor and inductor parser values. \\
18 & \path{18_diode_model_values} & Exact numeric parsing of diode model
parameters. \\
19 & \path{19_dependent_source_parser_bundle} & Parser bundle for E, G, F,
and H source types. \\
20 & \path{20_diode_parser_symbolic_model} & Symbolic diode model parsing
and isolated diode residual. \\
21 & \path{21_single_capacitor_frequency} & Isolated capacitor and
\(s=j2\pi f\) substitution. \\
22 & \path{22_numeric_substitution_missing_params} & Strict missing-parameter
error. \\
23 & \path{23_linear_extraction_rejects_diode} & Expected nonlinear rejection
from linear extraction. \\
24 & \path{24_descriptor_rejects_laplace} & Expected descriptor rejection in
Laplace domain. \\
25 & \path{25_case_preserved_component_name} & Case-insensitive component
type with preserved original name. \\
26 & \path{26_invalid_examples_expected_errors} & Expected parser and model
errors for invalid syntax. \\
\bottomrule
\end{longtable}
\endgroup

\section{Verification Commands}
\label{app:verification}

The example catalog can be checked from the repository root with:

\begin{lstlisting}[style=shell]
for f in examples/user_guide/*/run.py; do
  /home/konst/Documents/Projects/venv/bin/python "$f"
done
\end{lstlisting}

\begin{lstlisting}[style=shell]
for f in examples/user_guide/*/ngspice.cir; do
  ngspice -b "$f"
done
\end{lstlisting}

\begin{lstlisting}[style=shell]
/home/konst/Documents/Projects/venv/bin/ruff check examples/user_guide
\end{lstlisting}

\end{document}
